\documentclass[12pt,a4paper]{article}

% ── Encoding & Language ───────────────────────────────────────────
\usepackage[utf8]{inputenc}
\usepackage[T1]{fontenc}
\usepackage[spanish,english]{babel}
\usepackage{csquotes}

% ── Geometry & Spacing ────────────────────────────────────────────
\usepackage[top=3cm,bottom=2.5cm,left=2.5cm,right=2.5cm]{geometry}
\usepackage{setspace}
\onehalfspacing

% ── Graphics & Figures ────────────────────────────────────────────
\usepackage{graphicx}
\usepackage{tikz}
\usetikzlibrary{shapes.geometric, arrows.meta, positioning, fit, backgrounds}
\usepackage{float}

% ── Tables & Lists ────────────────────────────────────────────────
\usepackage{booktabs}
\usepackage{longtable}
\usepackage{array}
\usepackage{multirow}
\usepackage{enumitem}
\setlist{nosep,leftmargin=*}

% ── Maths ─────────────────────────────────────────────────────────
\usepackage{amsmath}
\usepackage{amssymb}
\usepackage{siunitx}
\sisetup{detect-all}

% ── Hyperlinks & Metadata ─────────────────────────────────────────
\usepackage[colorlinks=true,linkcolor=blue,citecolor=blue,urlcolor=blue]{hyperref}
\usepackage{url}

% ── Captions ──────────────────────────────────────────────────────
\usepackage{caption}
\usepackage{subcaption}

% ── Title Metadata ────────────────────────────────────────────────
\title{
  \vspace{-2cm}
  \textbf{Propuesta de Diseño Conceptual}\\[0.3cm]
  \large Sistema de Monitoreo SDR de Exposición a Campos Eléctricos\\
  en Entornos Interiores con Alta Densidad IoT
}
\author{
  \small
  Grupo de Control y Procesamiento Digital de Señales (GCPDS)\\
  Universidad Nacional de Colombia --- Sede Manizales\\
  \vspace{0.2cm}\\
  Universidad de Caldas
}
\date{\small \today}

\begin{document}

\maketitle
\thispagestyle{empty}

% ── Abstract ──────────────────────────────────────────────────────
\begin{abstract}
\noindent
Este documento presenta el diseño conceptual de una plataforma IoT basada en Radio Definida por Software (SDR) para el monitoreo de exposición humana a campos eléctricos en entornos interiores con alta densidad de dispositivos IoT. La arquitectura propuesta consta de cinco componentes principales: hardware SDR, motor de adquisición en tiempo real en C, procesamiento digital de señales y lógica de negocio en Python, base de datos temporal TimescaleDB, y una interfaz web reactiva. El sistema emplea técnicas de aprendizaje automático para umbralización adaptativa, sensado cooperativo de espectro, e interpolación geoestadística (Kriging) para la generación de mapas térmicos de exposición electromagnética, todo en conformidad con las normas ICNIRP y UIT-T K.52. Se incluye la lista de materiales (BOM) del hardware, la cadena de señal RF, y la planificación temporal del proyecto.
\end{abstract}

\vspace{1cm}
{\small
\noindent\textbf{Palabras clave:} SDR, IoT, exposición electromagnética, Kriging, aprendizaje automático, USRP B210, Jetson Orin, mapas térmicos.
}

\newpage
\tableofcontents
\newpage

% ═══════════════════════════════════════════════════════════════════════
% 1. INTRODUCTION
% ═══════════════════════════════════════════════════════════════════════
\section{Introducción}

\subsection{Planteamiento del Problema}

La proliferación de dispositivos IoT en entornos laborales cerrados (oficinas, fábricas, centros de datos) genera una preocupación creciente sobre la exposición acumulativa a campos eléctricos. Los analizadores de espectro tradicionales son costosos, voluminosos y poco prácticos para estudios densos o móviles en interiores. El Radio Definido por Software (SDR) surge como una alternativa de menor costo y mayor flexibilidad, aunque con limitaciones en rango dinámico, comportamiento ADC/DAC y banda de frecuencia utilizable.

El sistema debe operar en el rango propuesto de \qtyrange{700}{3500}{\MHz}, abarcando las principales bandas IoT y de telecomunicaciones (LTE, WiFi 2.4~GHz, ISM). Los entornos interiores requieren numerosas mediciones rápidas para identificar zonas de alta radiación y compararlas con las normas de exposición.

\subsection{Justificación}

La umbralización estática resulta insuficiente en entornos IoT ruidosos y dinámicos. La propuesta incorpora umbralización adaptativa basada en aprendizaje automático y sensado cooperativo. Los mapas de exposición de campo eléctrico permiten correlacionar densidad de dispositivos, uso frecuencial y niveles de exposición, constituyendo una herramienta tanto científica como regulatoria.

\subsection{Pregunta de Investigación}

\noindent
\textit{¿Cómo influye la densidad de dispositivos IoT en los niveles de exposición a campos eléctricos en entornos cerrados y qué impacto puede tener dicha exposición en la salud de los trabajadores, utilizando un sistema de monitoreo basado en SDR y técnicas avanzadas de aprendizaje automático para la detección y análisis?}

\subsection{Objetivos}

\paragraph{Objetivo General}
Desarrollar un sistema de evaluación de niveles de exposición humana a campos eléctricos, basado en SDR y técnicas de aprendizaje automático, que permita medir la densidad de radiación en entornos interiores con alta ocupación de dispositivos IoT.

\paragraph{Objetivos Específicos}
\begin{enumerate}[label=OE\arabic*:,leftmargin=2.5em]
  \item Diseñar y construir un módulo de medición SDR para monitoreo espectral de redes IoT de alta densidad en espacios interiores, usando técnicas de sensado de espectro.
  \item Desarrollar algoritmos de umbralización y procesos cooperativos basados en métodos de aprendizaje automático para manejar la variabilidad de mediciones en entornos interiores con alta interferencia.
  \item Generar mapas espaciales de campos eléctricos usando técnicas avanzadas de medición y análisis de datos para evaluar conformidad con regulaciones de exposición humana.
\end{enumerate}

% ═══════════════════════════════════════════════════════════════════════
% 2. SYSTEM ARCHITECTURE
% ═══════════════════════════════════════════════════════════════════════
\section{Arquitectura del Sistema}

\subsection{Visión General}

El sistema se compone de cinco módulos principales que operan en conjunto para capturar, procesar, almacenar y visualizar datos de exposición a campos eléctricos. La arquitectura está diseñada bajo los principios de procesamiento en tiempo real, baja latencia y desacoplamiento de componentes, permitiendo la evolución independiente de cada módulo.

\begin{figure}[H]
\centering
\begin{tikzpicture}[
  node distance=1.2cm and 1.5cm,
  box/.style={rectangle, draw=black, rounded corners, minimum width=4cm, minimum height=0.9cm, align=center, font=\small\sffamily},
  sbox/.style={box, minimum width=3.2cm, minimum height=0.7cm, font=\footnotesize\sffamily},
  >={Stealth[scale=0.8]},
  line/.style={draw, ->, >=Stealth},
  doubleline/.style={draw, <->, >=Stealth},
]
  % ── SDR ────────────────────────────────────
  \node[box, fill=red!15] (sdr) {SDR\\Hardware};

  % ── C Layer ────────────────────────────────
  \node[box, fill=blue!12, right=of sdr, minimum height=4.5cm, minimum width=6cm] (clayer) {};
  \node[font=\small\sffamily\bfseries, above=0.1cm of clayer.north] {Capa C (Tiempo Real)};
  \node[sbox, fill=blue!8, above=1.2cm of clayer.south] (buffers) {Buffers};
  \node[sbox, fill=blue!8, below=0.0cm of buffers] (engine) {Motor de Eventos\\(No Polling)};
  \node[sbox, fill=blue!8, below=0.0cm of engine] (psd) {PSD + Modelo IA};

  % ── Python Layer ───────────────────────────
  \node[box, fill=green!12, right=of clayer, minimum height=3cm, minimum width=4.5cm] (python) {};
  \node[font=\small\sffamily\bfseries, above=0.1cm of python.north] {Capa Python (DSP)};
  \node[sbox, fill=green!8, above=0.8cm of python.south] (dsp) {DSP-RF};
  \node[sbox, fill=green!8, below=0.0cm of dsp] (queries) {Consultas (Queries)};

  % ── Database ───────────────────────────────
  \node[box, fill=yellow!15, below=of python, minimum width=4cm] (db) {Base de Datos\\TimescaleDB (PostgreSQL)};

  % ── Frontend ───────────────────────────────
  \node[box, fill=purple!12, right=of python, minimum height=2.5cm, minimum width=4cm] (frontend) {};
  \node[font=\small\sffamily\bfseries, above=0.1cm of frontend.north] {Frontend};
  \node[sbox, fill=purple!8, above=0.5cm of frontend.south] (react) {Vite + React\\Mapas Térmicos, Espectro,\\ Dashboard, Reportes};

  % ── Connections ────────────────────────────
  \draw[doubleline] (sdr) -- node[above,font=\footnotesize] {RF Raw} (buffers);
  \draw[doubleline] (engine.west) -- ++(-0.5,0) |- node[pos=0.25,left,font=\footnotesize,align=center] {ZMQ\\IPC} (dsp.west);
  \draw[line, dashed] (dsp.west) -- ++(-0.3,0) |- node[pos=0.25,left,font=\footnotesize,align=center] {Medición\\Programada} (engine.west);
  \draw[line] (clayer.north east) -- ++(0.5,0) |- node[pos=0.25,right,font=\footnotesize,align=center] {WebSocket\\Binario} (frontend.north west);
  \draw[line, dashed] (frontend.south west) -- ++(-0.5,0) |- node[pos=0.25,right,font=\footnotesize,align=center] {Trigger,\\Parámetros} (clayer.south east);
  \draw[doubleline] (dsp.east) -- node[above,font=\footnotesize] {Query / Program} (react.west);
  \draw[doubleline] (queries.south) -- (db.north);

\end{tikzpicture}
\caption{Diagrama conceptual de la arquitectura de 5 componentes. Las líneas continuas representan flujos de datos primarios; las punteadas representan comandos de control.}
\label{fig:arch}
\end{figure}

\subsection{Componente 1: SDR (Radio Definida por Software)}

\textbf{Rol:} Dispositivo externo de captura RF.

El SDR es el punto de entrada de señales electromagnéticas al sistema. Captura muestras crudas de RF, cubriendo el rango de \qtyrange{698}{6000}{\MHz} con un ancho de banda instantáneo de hasta \qty{56}{\MHz}. El hardware seleccionado para el prototipo es el USRP B210 de Ettus Research, que ofrece arquitectura MIMO 2$\times$2 con resolución ADC/DAC de 12 bits e interfaz USB~3.0.

La conexión con el sistema es bidireccional: las muestras crudas fluyen hacia los buffers de la capa C, mientras que comandos de configuración (frecuencia, ganancia, tasa de muestreo) pueden ser enviados de vuelta al SDR.

\subsection{Componente 2: Capa C --- Motor de Tiempo Real}

\textbf{Rol:} Adquisición y preprocesamiento de señales de alto rendimiento.

La capa C es el núcleo de procesamiento en tiempo real del sistema. Está implementada en C11 y se ejecuta directamente sobre el hardware del host (no en contenedor), ya que requiere acceso directo a los periféricos USB del SDR. Se compone de tres submódulos:

\paragraph{Buffers}
Actúan como puente entre el hardware SDR y el procesamiento interno. Reciben muestras crudas del SDR en modo bidireccional y las entregan al motor de eventos, garantizando flujo continuo sin pérdida de datos.

\paragraph{Motor de Eventos (No Polling)}
Arquitectura dirigida por eventos que evita explícitamente el \textit{polling} para minimizar uso de CPU y latencia en el procesamiento de señales en tiempo real. Reacciona ante la llegada de datos desde los buffers y:
\begin{itemize}
  \item Envía datos y eventos a la capa Python mediante \textbf{ZMQ sobre IPC} (comunicación entre procesos) para procesamiento de alto nivel.
  \item Recibe comandos de control desde Python (\textit{medición programada}).
  \item Transmite datos espectrales binarios directamente al frontend web mediante \textbf{WebSocket} (Libwebsockets), evitando el paso por Python para flujos sensibles a latencia.
  \item Recibe comandos de control desde el frontend (\textit{trigger}, \textit{cambio de parámetros}).
\end{itemize}

\paragraph{PSD (Power Spectral Density) + Modelo IA}
Módulo de análisis de Densidad Espectral de Potencia embebido directamente en la capa C. Incorpora un modelo de inteligencia artificial para análisis espectral y detección de anomalías, operando cerca de la fuente de datos para inferencia de baja latencia. El modelo se entrena para umbralización adaptativa y reducción de falsas alarmas en entornos con alta interferencia.

\subsection{Componente 3: Capa Python --- DSP y Lógica de Negocio}

\textbf{Rol:} Procesamiento digital de señales de alto nivel, orquestación de datos y lógica de negocio.

La capa Python aprovecha el ecosistema científico (NumPy, SciPy) para procesamiento RF que no requiere la latencia mínima de la capa C. Se compone de dos submódulos:

\paragraph{DSP-RF}
Módulo de Procesamiento Digital de Señales enfocado en RF. Recibe eventos y datos desde la capa C vía ZMQ IPC, aplica algoritmos de análisis espectral avanzado, y envía comandos de medición programada de vuelta a la capa C. También se comunica bidireccionalmente con el frontend para operaciones de \textit{consulta} y \textit{programación} de campañas de medición.

\paragraph{Consultas (Queries)}
Capa de acceso a datos que gestiona toda la persistencia. Interactúa bidireccionalmente con la base de datos para almacenar mediciones, resultados de algoritmos, configuraciones, campañas y reportes. Implementa la lógica de negocio para trazabilidad científica: cada medición se relaciona con su configuración de adquisición, modelo de IA, campaña y ubicación espacial.

\subsection{Componente 4: Base de Datos --- TimescaleDB}

\textbf{Rol:} Almacenamiento persistente de series temporales con capacidades de consulta analítica.

TimescaleDB (extensión sobre PostgreSQL 16) proporciona:
\begin{itemize}
  \item Almacenamiento eficiente de series temporales (mediciones espectrales con timestamp).
  \item Compresión nativa y particionamiento automático (hypertables) para datos de alta frecuencia.
  \item Modelo relacional completo para metadatos: sensores, antenas, configuraciones, campañas, algoritmos, umbrales, alertas, usuarios y reportes de cumplimiento.
  \item Trazabilidad completa desde la captura cruda hasta el reporte exportable.
\end{itemize}

\subsection{Componente 5: Frontend --- Interfaz de Usuario}

\textbf{Rol:} Visualización reactiva para monitoreo, control y consulta.

Aplicación de página única (SPA) construida con Vite + React + TypeScript, diseñada con Material UI (tema oscuro). Los componentes principales de visualización incluyen:

\begin{itemize}
  \item \textbf{Espectro en Vivo:} Gráfico de cascada (waterfall) y FFT en tiempo real, alimentados directamente por el flujo binario WebSocket desde la capa C (Plotly.js).
  \item \textbf{Mapa Térmico IoT:} Mapa de exposición de campo eléctrico sobre planos de planta, generado mediante interpolación Kriging a partir de mediciones espaciales. Permite superposición de cuadrícula (\textit{grid}) y colocación de puntos de medición.
  \item \textbf{Estado de Sensores:} Monitoreo de salud, telemetría y conectividad de los sensores SDR desplegados.
  \item \textbf{Gestión de Campañas:} Programación de campañas de medición con asignación de sensores, rangos de frecuencia, resolución e intervalos.
  \item \textbf{Reportes Exportables:} Generación automática de reportes de cumplimiento normativo (ICNIRP, UIT-T K.52) con caché de resultados.
  \item \textbf{Panel de Configuración:} Estado de servicios, parámetros del sistema y control de acceso.
\end{itemize}

% ═══════════════════════════════════════════════════════════════════════
% 3. COMMUNICATION CHANNELS
% ═══════════════════════════════════════════════════════════════════════
\section{Canales de Comunicación}

La Tabla~\ref{tab:channels} resume los canales de comunicación entre componentes, el protocolo utilizado y la dirección del flujo de datos.

\begin{table}[H]
\centering
\caption{Canales de comunicación del sistema}
\label{tab:channels}
\begin{tabular}{@{}llll@{}}
\toprule
\textbf{Origen} & \textbf{Destino} & \textbf{Protocolo} & \textbf{Dirección} \\
\midrule
SDR & Buffers (C) & Muestras RF crudas & Bidireccional \\
Buffers (C) & Motor de Eventos (C) & Flujo interno & Interno \\
Motor de Eventos (C) & DSP-RF (Python) & ZMQ sobre IPC & Bidireccional \\
DSP-RF (Python) & Motor de Eventos (C) & Medición programada & Python $\to$ C \\
Capa C & Frontend (React) & WebSocket binario & C $\to$ Frontend \\
Frontend (React) & Capa C & WebSocket (control) & Frontend $\to$ C \\
DSP-RF (Python) & Frontend (React) & HTTP (REST) & Bidireccional \\
Consultas (Python) & Base de Datos & SQL & Bidireccional \\
\bottomrule
\end{tabular}
\end{table}

\subsection{Decisiones de Diseño de Comunicación}

\begin{description}
  \item[ZMQ sobre IPC] ZeroMQ con comunicación entre procesos (IPC) proporciona mensajería de baja latencia entre los procesos C y Python que coexisten en el mismo host, eliminando la sobrecarga de la pila de red.
  \item[WebSocket binario directo] El flujo de datos espectrales se transmite directamente desde la capa C al navegador sin pasar por Python. Esto minimiza la latencia para visualizaciones en tiempo real (cascada, FFT) que requieren alta tasa de actualización.
  \item[HTTP REST para consultas] Las operaciones de consulta y programación entre frontend y backend Python utilizan HTTP/REST, adecuado para operaciones transaccionales que no requieren flujo continuo.
\end{description}

% ═══════════════════════════════════════════════════════════════════════
% 4. HARDWARE BILL OF MATERIALS
% ═══════════════════════════════════════════════════════════════════════
\section{Lista de Materiales de Hardware (BOM)}

El sistema de monitoreo es portátil y autocontenido, diseñado para capturar, procesar y visualizar señales RF en el rango de \qtyrange{698}{6000}{\MHz}. A continuación se detalla cada componente de hardware con su rol, especificaciones, proveedor y costo.

\subsection{Resumen de Costos}

\begin{table}[H]
\centering
\caption{Resumen de costos del hardware}
\label{tab:cost}
\begin{tabular}{@{}lrr@{}}
\toprule
\textbf{Componente} & \textbf{Precio Unitario (USD)} & \textbf{Precio Importado (USD)} \\
\midrule
USRP B210 & \$2,200.00 & \$2,860.00 \\
Antena 698--2700 MHz & \$200.00   & \$260.00   \\
Filtros de Preselección ($\times$2) & \$200.00   & \$260.00   \\
Jetson Orin Nano Super & \$300.00   & \$390.00   \\
Pantalla Táctil 7'' & \$60.00    & \$78.00    \\
Carcaza Personalizada & \$100.00   & \$130.00   \\
\midrule
\textbf{TOTAL} & \textbf{\$3,060.00} & \textbf{\$3,978.00} \\
\bottomrule
\end{tabular}
\end{table}

\subsection{Especificaciones por Componente}

\subsubsection{SDR --- Ettus Research USRP B210}

\textbf{Rol:} Front-end RF principal. Captura de espectro de banda ancha y transmisión vía USB 3.0 al módulo de cómputo.

\begin{table}[H]
\centering
\caption{Especificaciones del USRP B210}
\begin{tabular}{@{}ll@{}}
\toprule
\textbf{Parámetro} & \textbf{Valor} \\
\midrule
Modelo / Parte & USRP B210 (Kit, UB210-KIT) \\
Proveedor & Ettus Research (National Instruments) \\
Rango de frecuencia & \qtyrange{70}{6000}{\MHz} continuo \\
Ancho de banda RF & Hasta \qty{56}{\MHz} instantáneo \\
Arquitectura & MIMO 2$\times$2 (2 TX, 2 RX) \\
ADC/DAC & 12 bits \\
Interfaz host & USB 3.0 \\
FPGA & Xilinx Spartan-6 XC6SLX150 \\
Figura de ruido & $\approx$\qty{8}{\dB} \\
Impedancia & \qty{50}{\ohm} \\
Conectores RF & 4 SMA (RX/TX por canal) \\
Software soportado & GNU Radio, UHD, MATLAB, LabVIEW \\
\bottomrule
\end{tabular}
\end{table}

\subsubsection{Módulo de Cómputo --- NVIDIA Jetson Orin Nano Super}

\textbf{Rol:} Unidad central de procesamiento para inferencia de IA, procesamiento de señales, control del sistema y salida de pantalla.

\begin{table}[H]
\centering
\caption{Especificaciones del Jetson Orin Nano Super}
\begin{tabular}{@{}ll@{}}
\toprule
\textbf{Parámetro} & \textbf{Valor} \\
\midrule
Modelo & Jetson Orin Nano Super Developer Kit (945-13766-0005-000) \\
Rendimiento IA & Hasta 67 TOPS \\
GPU & NVIDIA Ampere, 1024 núcleos CUDA \\
CPU & 6 núcleos ARM Cortex-A78AE \\
RAM & 8~GB LPDDR5, 102~GB/s ancho de banda \\
Almacenamiento & 2$\times$ M.2 Key M (NVMe SSD) \\
USB & 4$\times$ USB 3.2 Gen 2 Type-A, 1$\times$ USB-C \\
Video & DisplayPort 1.2 (HDMI vía adaptador) \\
Conectividad & Wi-Fi (M.2 Key E), Gigabit Ethernet \\
GPIO & Header de 40 pines \\
Sistema operativo & JetPack (Ubuntu Linux) \\
Soporte IA & TensorRT, PyTorch, ROS 2, GNU Radio \\
\bottomrule
\end{tabular}
\end{table}

\subsubsection{Antena Omnidireccional --- L-com LCANOM1075}

\textbf{Rol:} Elemento principal de captación de señal RF en el rango de \qtyrange{698}{2700}{\MHz}.

\begin{table}[H]
\centering
\caption{Especificaciones de la antena}
\begin{tabular}{@{}ll@{}}
\toprule
\textbf{Parámetro} & \textbf{Valor} \\
\midrule
Modelo & LCANOM1075 \\
Fabricante & L-com \\
Rango de frecuencia & \qtyrange{698}{2700}{\MHz} \\
Ganancia & 6.5~dBi \\
Patrón de radiación & Omnidireccional \\
Polarización & Vertical \\
Conector & SMA Macho (resorte, \qty{5.5}{\mm}) \\
Radomo & Fibra de vidrio, negro \\
Impedancia & \qty{50}{\ohm} \\
VSWR & $\leq$ 2.0 (típico) \\
\bottomrule
\end{tabular}
\end{table}

\subsubsection{Filtros de Preselección --- Mini-Circuits}

\textbf{Rol:} Filtrado pasabanda entre la antena y el USRP B210. Suprimen interferencia fuera de banda antes de la digitalización, mejorando la sensibilidad y reduciendo la intermodulación. Se utilizan dos filtros, uno por cada canal RX del USRP B210.

\begin{table}[H]
\centering
\caption{Filtro 1 --- VBFZ-925-S+}
\begin{tabular}{@{}ll@{}}
\toprule
\textbf{Parámetro} & \textbf{Valor} \\
\midrule
Modelo & VBFZ-925-S+ \\
Fabricante & Mini-Circuits \\
Tipo & Filtro pasabanda conectorizado \\
Pasabanda & \qtyrange{800}{1050}{\MHz} \\
Conector & SMA \\
Impedancia & \qty{50}{\ohm} \\
Cobertura espectral & GSM 850, GSM 900, LTE B5/B8/B20 \\
\bottomrule
\end{tabular}
\end{table}

\begin{table}[H]
\centering
\caption{Filtro 2 --- BPF-A1340+}
\begin{tabular}{@{}ll@{}}
\toprule
\textbf{Parámetro} & \textbf{Valor} \\
\midrule
Modelo & BPF-A1340+ \\
Fabricante & Mini-Circuits \\
Tipo & Filtro pasabanda LC \\
Pasabanda & \qtyrange{1000}{1800}{\MHz} \\
Conector & SMA \\
Impedancia & \qty{50}{\ohm} \\
Cobertura espectral & GPS L1 (1575 MHz), DCS 1800, LTE B1/B3/B4/B10 \\
\bottomrule
\end{tabular}
\end{table}

La combinación de ambos filtros proporciona preselectividad en el rango de \qtyrange{800}{1800}{\MHz}, permitiendo monitoreo simultáneo de dos bandas independientes.

\subsubsection{Pantalla Táctil --- Waveshare 7'' HDMI LCD (C)}

\textbf{Rol:} Interfaz de usuario táctil para visualización de espectro, mapas térmicos, dashboards y estado del sistema.

\begin{table}[H]
\centering
\caption{Especificaciones de la pantalla}
\begin{tabular}{@{}ll@{}}
\toprule
\textbf{Parámetro} & \textbf{Valor} \\
\midrule
Modelo & Waveshare 7'' HDMI LCD (C), Rev 4.1 \\
Tamaño & 7~pulgadas \\
Tipo de panel & IPS \\
Resolución & 1024 $\times$ 600 \\
Táctil & Capacitivo, 5 puntos \\
Interfaz de video & HDMI \\
Interfaz táctil & USB Micro-B (HID) \\
Sistemas soportados & Ubuntu, Raspbian, Windows, JetPack \\
\bottomrule
\end{tabular}
\end{table}

\subsubsection{Carcaza Personalizada}

\textbf{Rol:} Contenedor mecánico para todos los componentes. Proporciona protección física, gestión de cables, puntos de montaje para la antena y la pantalla.

\textbf{Requisitos de diseño:}
\begin{itemize}
  \item Alojar PCB del USRP B210, Jetson Orin Nano, pantalla de 7~pulgadas.
  \item Puerto de montaje externo para antena (pasamuros SMA).
  \item Acceso a puerto USB 3.0 para conexión del USRP B210.
  \item Ventilación para disipación térmica (Jetson $\approx$ \qty{15}{\W}).
  \item Ventana para pantalla ($\approx$ \qtyproduct{165 x 100}{\mm} de área visible).
  \item Factor de forma portátil, tipo maletín de instrumento robusto.
\end{itemize}

\subsection{Cadena de Señal RF}

La Figura~\ref{fig:signalchain} muestra el flujo de la señal RF desde la antena hasta la visualización.

\begin{figure}[H]
\centering
\begin{tikzpicture}[
  node distance=0.8cm,
  block/.style={rectangle, draw, rounded corners, minimum width=5.5cm, minimum height=0.8cm, align=center, font=\footnotesize\sffamily},
  >={Stealth[]},
]
  \node[block, fill=gray!10] (ant) {ANTENA\\L-com LCANOM1075 --- \qtyrange{698}{2700}{\MHz}};
  \node[block, fill=blue!10, below=of ant] (filtros) {FILTROS DE PRESELECCIÓN\\VBFZ-925-S+ (800--1050 MHz) + BPF-A1340+ (1000--1800 MHz)};
  \node[block, fill=red!10, below=of filtros] (usrp) {SDR --- USRP B210\\2$\times$2 MIMO, \qty{56}{\MHz} BW instantáneo};
  \node[block, fill=green!10, below=of usrp] (jetson) {CÓMPUTO --- Jetson Orin Nano Super\\67 TOPS, GNU Radio, IA (TensorRT / PyTorch)};
  \node[block, fill=purple!10, below=of jetson] (display) {VISUALIZACIÓN --- Waveshare 7''\\HDMI + USB Touch};
  \node[draw, dashed, rounded corners, fit=(ant)(jetson), inner sep=10pt, label={[font=\footnotesize\sffamily]below:Carcaza Personalizada}] {};

  \draw[->,thick] (ant) -- node[right,font=\footnotesize] {SMA} (filtros);
  \draw[->,thick] (filtros) -- node[right,font=\footnotesize] {SMA} (usrp);
  \draw[->,thick] (usrp) -- node[right,font=\footnotesize] {USB 3.0} (jetson);
  \draw[->,thick] (jetson) -- node[right,font=\footnotesize] {HDMI + USB} (display);
\end{tikzpicture}
\caption{Cadena de señal RF: desde la antena hasta la visualización.}
\label{fig:signalchain}
\end{figure}

% ═══════════════════════════════════════════════════════════════════════
% 5. AI AND INTERPOLATION STRATEGY
% ═══════════════════════════════════════════════════════════════════════
\section{Estrategia de IA e Interpolación}

\subsection{Umbralización Adaptativa y Sensado Cooperativo}

El sistema de umbralización estática tradicional no es adecuado para entornos IoT con alta densidad de dispositivos, donde los niveles de ruido e interferencia varían significativamente. La propuesta incorpora:

\begin{itemize}
  \item \textbf{Autoencoders} para extracción de características y reducción de dimensionalidad en datos espectrales crudos. Permiten identificar patrones de interferencia característicos de entornos IoT y diferenciarlos de anomalías que requieren atención.
  \item \textbf{Umbralización de dos niveles} inspirada en sistemas cooperativos de comunicación (como 5G): un primer umbral rápido para detección gruesa y un segundo umbral refinado, validado cooperativamente entre múltiples barridos o frecuencias, para reducción de falsas alarmas.
  \item \textbf{Modelo de IA embebido en la capa C} (módulo PSD) que opera con inferencia de baja latencia, cerca de la fuente de datos, sin requerir transferencia al backend Python para decisiones en tiempo real.
\end{itemize}

\subsection{Interpolación Kriging}

Para la generación de mapas espaciales de exposición a campos eléctricos se emplea interpolación Kriging, un método geoestadístico que:

\begin{itemize}
  \item Estima valores en ubicaciones no muestreadas a partir de un conjunto discreto de puntos de medición con coordenadas geográficas (latitud, longitud).
  \item Proporciona no solo la estimación puntual, sino también la varianza del error de estimación, permitiendo evaluar la confiabilidad del mapa generado.
  \item Reduce el número de puntos de medición necesarios mediante muestreo selectivo inteligente, preservando la calidad del mapa y optimizando el tiempo de caracterización de un entorno interior.
\end{itemize}

% ═══════════════════════════════════════════════════════════════════════
% 6. DATA MODEL CONCEPT
% ═══════════════════════════════════════════════════════════════════════
\section{Modelo de Datos Conceptual}

El modelo de datos está diseñado para garantizar trazabilidad científica completa: cada medición puede ser rastreada hasta su configuración de adquisición, sensor, campaña, modelo de IA y ubicación.

\begin{figure}[H]
\centering
\begin{tikzpicture}[
  node distance=1.2cm and 1.5cm,
  entity/.style={rectangle, draw, rounded corners, minimum width=3cm, minimum height=0.7cm, align=center, font=\small\sffamily},
  rel/.style={draw, ->, >=Stealth},
]
  \node[entity, fill=blue!10] (sensor) {Sensores\\+ Antenas};
  \node[entity, fill=blue!10, right=of sensor] (config) {Configuración\\de Adquisición};
  \node[entity, fill=green!10, below=of config] (data) {Mediciones\\(Sensor Data)};
  \node[entity, fill=yellow!10, below=of sensor] (campaign) {Campañas\\de Medición};
  \node[entity, fill=red!10, right=4.5cm of data] (algo) {Algoritmos\\+ Umbrales};
  \node[entity, fill=purple!10, below=of data] (report) {Reportes de\\Cumplimiento};
  \node[entity, fill=gray!10, right=of campaign] (user) {Usuarios\\+ Auditoría};

  \draw[rel] (sensor) -- (config);
  \draw[rel] (sensor) -- (data);
  \draw[rel] (config) -- (data);
  \draw[rel] (campaign) -- (data);
  \draw[rel] (algo) -- (data);
  \draw[rel] (data) -- (report);
  \draw[rel] (campaign) -- (report);
  \draw[rel] (user) -- (campaign);
\end{tikzpicture}
\caption{Modelo conceptual de datos: trazabilidad desde el sensor hasta el reporte de cumplimiento.}
\label{fig:datamodel}
\end{figure}

Las entidades principales del modelo incluyen: sensores, antenas, configuraciones de adquisición, datos de medición (con coordenadas espaciales), campañas, algoritmos (con versiones), umbrales de exposición configurables, reportes de cumplimiento y registros de auditoría.

% ═══════════════════════════════════════════════════════════════════════
% 7. TIMELINE
% ═══════════════════════════════════════════════════════════════════════
\section{Planificación Temporal}

\subsection{Cronograma de Ejecución (12 Semanas)}

El plan de ejecución práctico se organiza en tres meses con hitos definidos:

\begin{table}[H]
\centering
\caption{Cronograma de 12 semanas}
\label{tab:schedule}
\begin{tabular}{@{}p{2.5cm}p{2.5cm}p{9cm}@{}}
\toprule
\textbf{Mes} & \textbf{Semanas} & \textbf{Actividades Principales} \\
\midrule
\multirow{3}{*}{Mes 1} & 1 & Escritura de artículo; auditoría de base de datos; documentación matemática de algoritmos de adquisición; configuración inicial del hardware SDR. \\
& 2 & Scripts de limpieza de datos históricos; definición de parámetros de umbralización para detección IoT; validación de tasa de transferencia sensor--base de datos. \\
& 3 & Conexión API / dashboard; entorno Autoencoder para extracción de características; ajustes de antenas y posicionamiento de dispositivos. \\
& 4 & Primeras pruebas de visualización web con puntos de datos aislados; pruebas de filtrado con datos reales; estabilidad operacional SDR (TRL 7). \\
\midrule
\multirow{3}{*}{Mes 2} & 5 & Diseño UI de mapa térmico para oficina; implementación de interpolación Kriging; recolección sistemática de datos. \\
& 6 & Entrenamiento de Autoencoder con datos limpios; integración de lógica de algoritmos en plataforma visual; soporte a calibración de sensores. \\
& 7 & Optimización del dashboard para renderizado en tiempo real; ajuste de modelo matemático según normas ICNIRP; refinamiento de IA para reducir falsas alarmas. \\
& 8 & Prueba integrada completa: captura $\to$ procesamiento/IA $\to$ visualización. \\
\midrule
\multirow{3}{*}{Mes 3} & 9 & Validación de precisión del mapa frente a mediciones de control manual; reportes exportables automáticos; documentación de lógica cooperativa. \\
& 10 & Desmontaje o fijación final del prototipo; validación de rendimiento del Autoencoder; corrección de bugs y seguridad de acceso. \\
& 11 & Pruebas finales del sistema; entrega de manual de usuario y manual técnico. \\
& 12 & Cierre de plan de trabajo, revisión de hitos y preparación de presentación de resultados. \\
\bottomrule
\end{tabular}
\end{table}

\subsection{Hitos Principales}

\begin{enumerate}[leftmargin=*]
  \item \textbf{Semana 4:} Datos de sensor fluyendo a la plataforma de forma organizada.
  \item \textbf{Semana 8:} Primera versión funcional de mapa térmico con interpolación Kriging.
  \item \textbf{Semana 12:} Sistema operativo, validado y con entregables finales.
\end{enumerate}

% ═══════════════════════════════════════════════════════════════════════
% 8. REGULATORY COMPLIANCE
% ═══════════════════════════════════════════════════════════════════════
\section{Cumplimiento Normativo}

El sistema está diseñado para evaluar exposición electromagnética conforme a:

\begin{itemize}
  \item \textbf{ICNIRP} (Comisión Internacional de Protección contra Radiación No Ionizante): Directrices para limitación de exposición a campos eléctricos y magnéticos.
  \item \textbf{UIT-T K.52}: Orientación sobre el cumplimiento de los límites de exposición de las personas a los campos electromagnéticos.
  \item \textbf{Regulación de espectro en Colombia} por la ANE (Agencia Nacional del Espectro).
\end{itemize}

Los umbrales de exposición se modelan como datos de dominio configurables (no valores fijos en código), permitiendo actualización cuando las normas evolucionen.

% ═══════════════════════════════════════════════════════════════════════
% 9. EXPECTED RESULTS
% ═══════════════════════════════════════════════════════════════════════
\section{Resultados Esperados}

\begin{itemize}
  \item Sistema funcional de monitoreo de exposición a campos eléctricos validado en entornos reales con nivel TRL 7.
  \item Prototipo integrado de hardware y software con arquitectura de 5 componentes.
  \item Algoritmos de umbralización adaptativa y sensado cooperativo con reducción de falsas alarmas.
  \item Generación de mapas espaciales de exposición mediante Kriging.
  \item Al menos una publicación científica indexada en revista categorizada por Minciencias.
  \item Registro de software ante la Dirección Nacional de Derecho de Autor.
  \item Repositorio público con código fuente, instaladores, manuales y documentación de soporte.
  \item Reportes técnicos y documentos de tesis de estudiantes vinculados.
\end{itemize}

% ═══════════════════════════════════════════════════════════════════════
% PARTICIPANTS
% ═══════════════════════════════════════════════════════════════════════
\section{Participantes}

\begin{table}[H]
\centering
\caption{Equipo del proyecto}
\begin{tabular}{@{}lll@{}}
\toprule
\textbf{Rol} & \textbf{Nombre} & \textbf{Responsabilidad} \\
\midrule
Grupo de investigación & GCPDS, UNAL Manizales & Coordinación, cronograma, integración \\
Investigador principal & Julio César García Álvarez & Liderazgo científico y técnico \\
Coinvestigador & Marcelo Herrera González & Supervisión de análisis de información \\
Estudiante de pregrado & Alejandro Patiño Bedoya & Recolección de datos de campo \\
Estudiante de pregrado & Luis Felipe Giraldo Díaz & Reportes técnicos y artículos \\
Estudiante de maestría & Julián Andrés Salazar Parias & Diseño y optimización de algoritmos IA \\
\bottomrule
\end{tabular}
\end{table}

Entidades de apoyo: Universidad Nacional de Colombia, Universidad de Caldas, Open Business Consulting S.A.S., Inference S.A.S., Dunderlab S.A.S.

% ═══════════════════════════════════════════════════════════════════════
% REFERENCES
% ═══════════════════════════════════════════════════════════════════════
\begin{thebibliography}{9}

\bibitem{icnirp}
ICNIRP,
\textit{Guidelines for Limiting Exposure to Electromagnetic Fields (100~kHz to 300~GHz)},
Health Physics, vol.~118, no.~5, pp.~483--524, 2020.

\bibitem{uitk52}
ITU-T Recommendation K.52,
\textit{Guidance on complying with limits for human exposure to electromagnetic fields},
International Telecommunication Union, 2021.

\bibitem{kriging}
M.~A.~Oliver, R.~Webster,
\textit{Kriging: a method of interpolation for geographical information systems},
International Journal of Geographical Information Systems, vol.~4, no.~3, pp.~313--332, 1990.

\bibitem{usrp}
Ettus Research,
\textit{USRP B210 Product Overview},
\url{https://www.ettus.com/all-products/ub210-kit/}

\bibitem{jetson}
NVIDIA Corporation,
\textit{Jetson Orin Nano Super Developer Kit},
\url{https://developer.nvidia.com/embedded/jetson-orin-nano}

\bibitem{gnuradio}
GNU Radio Project,
\textit{GNU Radio: The Free and Open Software Radio Ecosystem},
\url{https://www.gnuradio.org/}

\bibitem{timescaledb}
Timescale Inc.,
\textit{TimescaleDB Documentation},
\url{https://docs.timescale.com/}

\end{thebibliography}

\end{document}
